\chapter{Rendszerszintű értékelés és megvitatás}

\section{A legfontosabb rendszertervezési döntések}

Az egész projektet visszanézve a legerősebb döntésnek az tekinthető, hogy a különböző feladatok nem kerültek egyetlen univerzális technológiába kényszerítve. A rendszer nem eszközpreferencia, hanem minőségi attribútumok mentén fejlődött: reprodukálhatóság, magyarázhatóság, perzisztencia, futási teljesítmény és demo-koherencia.

A Python gyors iterációt és kényelmes adatfeldolgozást ad, a Rust teljesítményt és pontos gráfkontrollt, a Node/Express stabil integrációs felületet, a React pedig erős bemutathatóságot.

\section{A rendszer erősségei}

A projekt legfontosabb erőssége, hogy a teljes adatút végigkövethető a Riot API-tól a frontendig. A játékosprofilok explicit, ellenőrizhető képleteken és mezőkön alapulnak; a gráfépítés és klaszterezés nem csak vizuális export, hanem perzisztens objektum, amelynek bejárása és újraelőállítása determinisztikus. A Rust réteg valódi analitikai és futtatási értéket ad. Az asszortativitás, a core-periphery gráfolvasat és a Brandes-centralitás valódi kutatási kérdésekké emelik a rendszert.

\section{Érvényességi fenyegetések}

A projekt legfontosabb validitási kockázata, hogy a dataset nem reprezentálja a teljes játékospopulációt: a gyűjtés seed- és queue-függő. Az \textit{opscore} és \textit{feedscore} értelmezhető, de nem objektív igazságmérő.

Rétegenként szétbontva: a \textbf{mintavételi validitás} szempontjából a seedválasztás, a premade-probe és a queue-orientáció együtt nem véletlen mintát adnak. A \textbf{mérési validitás} szempontjából az \textit{opscore} és \textit{feedscore} értelmezhetők, de konstrukciós döntésekre épülnek. A \textbf{projekciós validitás} szempontjából a gráfmód, az élküszöb és az ally/enemy kapcsolatértelmezés érdemben alakítja az analitikai mintát. A \textbf{klasztervaliditás} szempontjából a közösségdetektálás eredménye függ a választott gráftól és a modularitás felbontási korlátjától \cite{fortunato_barthelemy2007}. A \textbf{temporális validitás} szempontjából a korpusz fokozatosan épülő archívum.

Ezek a fenyegetések kijelölik azt a határt, amelyen belül az állítások védhetők.

\section{A megvalósult és a jövőbeli elemek szétválasztása}

A ténylegesen megvalósult mag: Riot API-ra épülő adatgyűjtés, \textit{puuid}-alapú normalizálás, \textit{feedscore}/\textit{opscore} számítás, SQLite klaszterperzisztencia, Rust Graph V2 gráfépítés, útvonalkeresés, asszortativitás, Brandes-centralitás és 3D globális export. Ehhez adódik a databaseproject és C\# böngésző alprojekt.

A temporal consistency és community cohesion vs. performance vizsgálatok eltávolítva a jelenlegi scope-ból.
